\subsection{Earlier bounds and related sunflower problems}\label{sec:history}

Erd\H{o}s and Rado suggested replacing the factorial in their 1960
bound by \(c^k\) for an absolute \(c\)~\cite[p.~86]{ER}. In the
notation of Theorem~\ref{thm:main}, this would give \(b_r=O(r)\).
Erd\H{o}s later stated the fixed-petal conjecture explicitly, allowing
the base to depend arbitrarily on \(r\)~\cite[Section~4]{ERstrong}.
For three petals the block-transversal lower bound can be improved.
A construction of Abbott, Hanson and Sauer gives
\(b_3\ge\sqrt{10}\)~\cite{AHS}.

After earlier improvements by Kostochka~\cite{Kostochka}, Alweiss,
Lovett, Wu and Zhang brought the base down to a polylogarithmic
function of \(k\) for fixed \(r\)~\cite{ALWZ}. Their argument finds a
subfamily whose members have a common intersection \(C\), such that
a random subset of the ground set is likely to contain one of the
nonempty remainders \(A\setminus C\). If a random color class in
an \(r\)-coloring contains no remainder with probability less than
\(1/r\), some coloring puts one remainder in each class. These
\(r\) disjoint remainders give the petals. Their examples also show why a logarithmic loss
arises here. For fixed sampling and success probabilities there are
families of size \((\log k)^{(1-o(1))k}\) with no such robust
sunflower~\cite[Section~3]{ALWZ}. Rao's coding argument gave
\(M_r(k)\le (Cr\log(rk))^k\)~\cite{Rao}. Bell, Chueluecha and Warnke
refined the probabilistic step to give \((Cr\log(k+1))^k\)~\cite{BCW}.
Related random-containment estimates enter the study of thresholds of
increasing families~\cite{Thresholds,PP}. Sunflower lemmas also play
a role in Razborov's lower bounds for monotone
circuits~\cite{Razborov}. Rao surveys these connections~\cite{RaoSurvey}.

For \(k\)-uniform families of VC dimension \(d\ge1\), Ge, Wang, Xu and Zhao
obtained the width-independent base \(50dr\)~\cite[Theorem~1.1]{VC},
building on earlier results for restricted set
systems~\cite{FPS,vcsmall}. Shangguan improved the base to
\(O(r\log(d+1))\)~\cite[Corollary~1.4]{Shangguan}. For unrestricted
families Fukuyama proposed a base \(O(r^4)\)~\cite{Fukuyama}.
A product-family obstruction to the sequence-selection estimate in
that argument is described in~\cite{Plasma}.

The size \(N\) of the ground set gives a different measure of the
problem. Erd\H{o}s and Szemer\'edi asked for a bound exponentially
smaller than \(2^N\) on a sunflower-free family of subsets of
an \(N\)-element set~\cite{ES}. Alon, Shpilka and Umans showed how
both conjectures can be studied using words, and that the fixed-petal
conjecture rules out a group construction proposed by Coppersmith
and Winograd for matrix multiplication~\cite{ASU}. Naslund and
Sawin gave a direct slice-rank proof of the ground-set bound for
three petals~\cite{NS}. Buci\'c, He and Lee obtained
\(3(2-1/(16r^3))^N\) for every \(r\ge3\) by polynomial
interpolation~\cite[Theorem~1]{BHL}. For words of length \(k\) over
\(D\) symbols, their bound has base
\((2r)^{2/(r-1)}D^{1-1/(r-1)}\)~\cite[Theorem~2]{BHL}. The
word estimate proved here has a single base for all alphabets, as
required for the passage to \(k\)-element sets on an arbitrary ground
set~\cite[Theorem~2.6]{ASU}.
